\textbf{а)}
\Proof
Пусть $f$ приводим над $\mathbb{Q}[x]$, то есть $f = gh$, причем $g, h \in \mathbb{Q}[x]$. Тогда $g = \frac{a}{b}g_1, h = \frac{c}{d}h_1$, где $g_1, h_1$ примитивные многочлены степени $\geq 1$ (и из $Z[x]$). Из билета 35 произведение примитивнх многочленов -- примитивный многочлен.

Тогда $f = \frac{ac}{bd}(g_1h_1)$, и (т.к $f, g_1h_1$ примитивны), то $\frac{ac}{bd} = 1$. Тогда $f$ приводим над $\mathbb{Z}[x]$, противоречие.
\EndProof

\textbf{б)}
\Proof
$f(x)$ примитивен, неприводим в $\mathbb{Z}[x]$, значит, неприводим в $\mathbb{Q}[x]$. В $\mathbb{Q}[x]$ это эквивалентно простоте. 

Пусть $f(x)|(g(x)h(x)), g, h \in \mathbb{Z}[x]$. Тогда в $\mathbb{Q}[x]$ один из них делится на $f$, пусть это $g(x)$ (без ограничения общности). Тогда $g(x) = f(x)a(x)$. Заметим, что $g(x) = bg_1(x), g_1(x)$ -- примитивный (в $\mathbb{Z}[x]$). А еще $a(x) = \frac{c}{d}a_1(x), a_1(x)$ -- примитивный (в $\mathbb{Z}[x]$).

Значит, $bdg_1(x) = ca_1(x)f(x)$. Произведение $a_1$ и $f$ -- примитивный многочлен. Тогда $bd = c$, и $g_1(x) = a_1(x)f(x)$. Значит, $f|g_1$, и $f|g$ в $\mathbb{Z}[x]$
\EndProof